\documentclass[12pt,a4paper]{article}

% --- Paquetes Esenciales ---
\usepackage[utf8]{inputenc}
\usepackage[T1]{fontenc}
\usepackage[spanish]{babel}
\usepackage[a4paper, margin=2.5cm]{geometry}
\usepackage{hyperref}
\usepackage{xcolor}

\hypersetup{
    colorlinks=true,
    linkcolor=blue,
    urlcolor=cyan,
}

\title{Propuestas de Proyectos de Robótica}
\author{Para: Sara \\ \small Preparado por el prof. César Rodríguez (Física-UDO) \\ \small con apoyo de Gemini Code Assist}
\date{\today}

\begin{document}

\maketitle

\section{Introducción}
La robótica es una forma increíble de aprender a resolver problemas uniendo la programación, la electrónica y la mecánica. A continuación, se presentan cinco ideas de proyectos pensados para ser accesibles, divertidos y fáciles de realizar en casa. La mayoría se pueden construir con materiales reciclados y un "Kit de Inicio" básico de plataformas como Arduino.

\section{Ideas de Proyectos}

\subsection{1. El Basurero Inteligente (Smart Trash Can)}
\textbf{Descripción:} Un basurero que abre su tapa automáticamente cuando detecta que una mano o un objeto se acerca, y se cierra unos segundos después. Es muy útil e higiénico. \\
\textbf{Materiales principales:} Un bote de basura pequeño con tapa abatible, una placa Arduino (o similar), un sensor ultrasónico (para medir la distancia) y un pequeño servomotor (para levantar la tapa mediante un hilo o alambre). \\
\textbf{Dificultad:} Baja. Es ideal para aprender cómo interactúan los sensores de distancia y el control de movimientos simples.

\subsection{2. Robot Evasor de Obstáculos}
\textbf{Descripción:} Un pequeño vehículo de dos ruedas que se desplaza de forma autónoma por la habitación. Cuando detecta un objeto frente a él (como una pared o un zapato), frena, gira hacia otro lado y sigue buscando un camino despejado. \\
\textbf{Materiales principales:} Un chasis de carrito para armar (muy económicos en línea), una placa Arduino, un controlador de motores (driver), un sensor ultrasónico y un portapilas. \\
\textbf{Dificultad:} Media. Es el "Hola Mundo" de los robots móviles. Enseña la lógica de las decisiones condicionales (si hay obstáculo a menos de 10 cm, entonces gira a la derecha).

\subsection{3. Brazo Robótico de Cartón (Mecánico o Electrónico)}
\textbf{Descripción:} La construcción de la estructura de un brazo robótico utilizando cartón duro, palillos y pegamento. Puede tener dos versiones:
\begin{itemize}
    \item \textit{Versión Hidráulica:} Controlado mediante jeringas y tubos de plástico con agua. ¡No requiere electrónica y es pura mecánica!
    \item \textit{Versión Electrónica:} Controlado mediante pequeños servomotores que se mueven al girar perillas (potenciómetros).
\end{itemize}
\textbf{Materiales principales:} Cartón, silicón, jeringas/tubos o, en su defecto, placa base, servomotores y potenciómetros. \\
\textbf{Dificultad:} Baja a Media (requiere paciencia manual). Desarrolla mucho la habilidad para armar estructuras y comprender cómo se transmiten los movimientos (cinemática).

\subsection{4. Sistema de Riego Automático}
\textbf{Descripción:} Un sistema que monitorea qué tan seca está la tierra de una planta y activa una pequeña bomba de agua de manera automática para regarla solo cuando tiene sed. \\
\textbf{Materiales principales:} Placa controladora, sensor de humedad de suelo, un módulo relé, una pequeña bomba de agua sumergible de 5V y un vaso con agua. \\
\textbf{Dificultad:} Baja. Aunque el movimiento no es un "vehículo", usa actuadores (la bomba) y sensores, combinando robótica con el cuidado del medio ambiente. Es excelente para ferias de ciencia.

\subsection{5. Girasol Robótico (Seguidor de Luz)}
\textbf{Descripción:} Una estructura (puede tener forma de flor o un pequeño panel) que utiliza sensores de luz para moverse siguiendo la fuente de luz más brillante en la habitación, tal como lo hace un girasol. \\
\textbf{Materiales principales:} Placa controladora, 2 fotorresistencias (LDR), un servomotor y material reciclado para el cuerpo. \\
\textbf{Dificultad:} Baja. Es un proyecto fascinante que muestra cómo la robótica a menudo se inspira en el comportamiento de la naturaleza.

\section{Consejos Finales para la Implementación de Proyectos de Robótica}
\begin{itemize}
    \item \textbf{Carcasa y Diseño:} ¡El diseño es tan importante como los cables! Decorar el proyecto con cartón pintado, materiales reciclados, figuras o piezas de Lego le da un toque único y personal que siempre destaca.
    \item \textbf{Paso a Paso:} Es mejor probar cada componente por separado. Por ejemplo, primero asegurarse de que el sensor pueda medir distancias y mostrarlo en la computadora, y luego intentar que ese sensor mueva el motor.
    \item \textbf{Uso de Kits:} Si deciden usar electrónica, adquirir un "Starter Kit" de Arduino es una excelente inversión, ya que trae guías, manuales y casi todos los componentes (sensores de luz, distancia, servomotores, cables y la placa) necesarios para hacer la mayoría de estos proyectos.
\end{itemize}

\end{document}